
%%% ====================================================================
%%% ====================================================================
%%%
%%%  SECTION 5: CATEGORICAL IMPLICATIONS AND FORMAL PROOF
%%%
%%% ====================================================================
%%% ====================================================================

\section{Categorical Implications and Formal Proof}\label{sec:categorical}

With the operator constructed (\cref{sec:methods}) and its predictions
verified against known zeros (\cref{sec:results}), we now interpret the
results within the categorical framework established by the Lean
formalization.  The central insight is that $\zeta(s)$ admits two
canonical decompositions---the Dirichlet series $\sum n^{-s}$
(additive, $\Set$-level, $\mu_2=1$) and the Euler product
$\prod_p(1-p^{-s})^{-1}$ (multiplicative, $\Met$-level,
$\mu_3=1/2$)---that factor through distinct levels of the
forgetful functor chain, and the norm contraction between those
levels determines the critical line.

\begin{remark}[The PCF framework as a diagram of functors]%
\label{rmk:pcf-diagram}
The PCF categorical framework is \emph{not} a single category but a
\textbf{diagram of categories connected by forgetful functors}:
\[
  C^*\text{-}\mathbf{Alg}
  \xrightarrow{\ U_1\ }
  \Hilb
  \xrightarrow{\ U_2\ }
  \Met
  \xrightarrow{\ U_3\ }
  \Set.
\]
The \emph{arity transition} (\cref{thm:arity-transition}) is a
property of this diagram: the forgetful functors contract the norm from
$\mu_2=1$ (binary, $\Set$ level) to $\mu_3=1/2$ (ternary, $\Met$ level).
Each category in the diagram is a standard mathematical category
($\Set$, $\Hilb$, $\Met$) which admits
\texttt{MonoidalCategory} structure; the coherence axioms (associator,
unitors, pentagon, triangle) are inherited from Mathlib's typeclass
system.  The ``PCF object'' $\Omega = P \cdot C \cdot F$ with
$\norm{\Omega}=1/2$ lives at the $\Met$ level, and the Lawvere
obstruction operates on the induced norm at each level.
\end{remark}

% ------------------------------------------------------------------
\subsection{\texorpdfstring{$\zeta$}{zeta} as a Categorical Object}\label{ssec:zeta-cat}
% ------------------------------------------------------------------

The Dirichlet series $\zeta(s)=\sum n^{-s}$ is additive---a sum over all
natural numbers---and factors through the binary level ($n=2$, $\mu_2=1$),
where the diagonal is permitted and Lawvere's theorem applies in full force.

The Euler product $\zeta(s)=\prod_p(1-p^{-s})^{-1}$ is multiplicative---a
product over primes.  Every prime $p>5$ satisfies
$p\bmod 20\in\Ztwenty$, where multiplication is closed and addition is
destroyed (\cref{prop:add-destroyed}).  The Euler product
therefore factors through the ternary level ($n=3$, $\mu_3=1/2$), where the
diagonal is blocked.

\begin{definition}[Golden residue group]\label{def:ZtwentyStar}
\begin{equation}\label{eq:ZtwentyStar}
  \Ztwenty = \{1,3,7,9,11,13,17,19\} \subset \mathbb{Z}/20\mathbb{Z}.
\end{equation}
\end{definition}

\begin{proposition}[Cardinality]\label{prop:ZtwentyStar-card}
\begin{equation}\label{eq:ZtwentyStar-card}
  |\Ztwenty| = 8.
\end{equation}
\end{proposition}

\begin{proof}
By enumeration: $\{1,3,7,9,11,13,17,19\}$ are exactly the integers
in $\{1,\ldots,19\}$ coprime to~$20$.  $|\cdot|=8=\varphi(20)$.
Verified by \lean{decide} in Lean.
\end{proof}

\begin{proposition}[Multiplicative closure and additive destruction]%
  \label{prop:add-destroyed}
The group $\Ztwenty$ is closed under multiplication in
$\mathbb{Z}/20\mathbb{Z}$:
\begin{equation}\label{eq:mul-preserved}
  a,b\in\Ztwenty \implies a\cdot b \in \Ztwenty,
\end{equation}
but \emph{not} under addition:
\begin{equation}\label{eq:addition-destroyed}
  \exists\, a,b\in\Ztwenty:\; a+b \notin \Ztwenty.
\end{equation}
\end{proposition}

\begin{proof}
Multiplication: $\Ztwenty$ is the group of units of
$\mathbb{Z}/20\mathbb{Z}$; closure under multiplication is the group
axiom.  The formal integrity of this closure is established
exhaustively for all $8\times 8 = 64$ pairs by \lean{decide} in Lean.

Addition: every element of $\Ztwenty$ is odd.  The sum of two odd
numbers is even, hence divisible by~$2$, hence not coprime to
$20=4\times 5$, hence not in~$\Ztwenty$.  Concretely, $1+3=4\notin\Ztwenty$.
The failure of additive closure is likewise verified exhaustively
for all 64 pairs by \lean{decide}. \end{proof}

\begin{proposition}[Exponent two]\label{prop:exponent-two}
The Galois group $G\cong\mathbb{Z}_2\times\mathbb{Z}_2$ satisfies:
\begin{equation}\label{eq:exponent-two}
  \forall g\in G,\quad g + g = 0.
\end{equation}
\end{proposition}

\begin{proof}
$\mathbb{Z}_2\times\mathbb{Z}_2 = \{(0,0),(1,0),(0,1),(1,1)\}$.
Each element $g$ satisfies $2g = 0$ since $2\equiv 0$ in $\mathbb{Z}_2$.
Verified by \lean{decide} (4 elements).
\end{proof}

\begin{proposition}[Fiber partition]\label{prop:fiber-partition}
The Galois homomorphism $G\colon\Ztwenty\to\mathbb{Z}_2\times\mathbb{Z}_2$
has four fibers that partition $\Ztwenty$ into four pairs.
\end{proposition}

\begin{proof}
$G$ is surjective onto $|\mathbb{Z}_2\times\mathbb{Z}_2|=4$ elements,
and $|\Ztwenty|=8$, so each fiber has exactly $8/4=2$ elements.  The
fibers are as listed in \cref{thm:golden-prime-classes}:
$\{1,9\}$, $\{11,19\}$, $\{13,17\}$, $\{3,7\}$.
\end{proof}

\begin{theorem}[Primes force ternary]\label{thm:zeta-forced}
\begin{equation}\label{eq:primes-in-ZtwentyStar}
  \forall p\text{ prime},\; p > 5 \implies
    p\bmod 20 \in \Ztwenty.
\end{equation}
\end{theorem}

\begin{proof}
A prime $p>5$ satisfies $\gcd(p,20)=1$ (since $20=2^2\cdot 5$ and
$p\nmid 2$, $p\nmid 5$), so $p\bmod 20\in\Ztwenty$, identifying the residue as a
unit in the group $\Ztwenty$.  In Lean: verified by \lean{fin_cases} over
$\mathbb{Z}/20\mathbb{Z}$, eliminating all non-unit residues by
divisibility.
\end{proof}

\begin{definition}[Arithmetic fragment]\label{def:arith-fragment}
The arithmetic fragment characterizes the theory through the distinction 
between two categorical arities~$n$:
\begin{equation}\label{eq:arith-fragment}
  \mathrm{ArithFragment} :=
    \mathrm{binary}~(n{=}2) \mid \mathrm{ternary}~(n{=}3).
\end{equation}
The unit group $\Ztwenty$ associated with the vigesimal period is ternary:
\begin{equation}\label{eq:ZtwentyStar-fragment}
  \mathrm{fragment}(\Ztwenty) = \mathrm{ternary}
  \qquad (|\Ztwenty| = 2^3).
\end{equation}
The cardinality $|\Ztwenty|=8$ identifies the group's combinatorial structure with the third-order hypercube $H_3$, the skeleton of the $n=3$ arity.
\end{definition}

\begin{proposition}[Fragment modulus]\label{prop:fragment-mu}
The spectral parameter $\mu$ is uniquely determined by the fragment's arity:
\begin{equation}\label{eq:fragment-mu-ternary}
  \mathrm{fragment\_mu}(\mathrm{ternary}) = \mu_3 = \tfrac{1}{2}.
\end{equation}
\end{proposition}

\begin{proof}
By \cref{eq:mu-n}, the ternary arity ($n=3$) yields $\mu_3 = 2 - 3/2 = 1/2$. This specific fragment modulus corresponds to the $\Met$ level where the diagonal is blocked (\cref{thm:arity-unique}).
\end{proof}

\begin{remark}
The Dirichlet series factors through the binary level ($\mu_2=1$);
the Euler product factors through the ternary level ($\mu_3=1/2$)
by \cref{thm:zeta-forced}.
\end{remark}


% ------------------------------------------------------------------
\subsection{PCF as a Contractive Monoidal Category}\label{ssec:CM-cat}
% ------------------------------------------------------------------

We formalize the norm obstruction in the language of monoidal categories.
The key construction is a \emph{tripartite object}~$\Omega$ whose norm
$\norm{\Omega}<1$ prevents the diagonal morphism required for Lawvere's
fixed-point theorem, while the specific value $\norm{\Omega}=1/2$ connects
to the spectral parameter~$\mu_3$.

%%% --- Spectral parametrization (CANONICAL DEFINITIONS) ---

\begin{definition}[Spectral parameters]\label{def:spectral-params}%
  \label{ssec:spectral-params}
For arity~$n$:
\begin{align}
  \sigma_n &= \frac{n}{2},
    \label{eq:sigma-n}  \\
  \mu_n &= 2 - \frac{n}{2}.
    \label{eq:mu-n}
\end{align}
These satisfy the spectral completeness relation:
\begin{equation}\label{eq:spectral-completeness}
  \sigma_n + \mu_n = 2 \qquad\text{for all } n.
\end{equation}
\end{definition}

\begin{theorem}[Arity uniqueness]\label{thm:arity-unique}
The conditions $\mu_n < 1$ (obstruction effective) and $\mu_n > 0$
(non-degenerate) hold simultaneously if and only if $n=3$.
\end{theorem}

\begin{proof}
$\mu_n < 1 \iff n > 2$; $\mu_n > 0 \iff n < 4$.
Together: $n=3$. \end{proof}

\begin{corollary}[Canonical spectral values]\label{cor:spectral-values}
At $n=3$:
\begin{align}
  \mu_2 &= 1 \qquad\text{(binary)},
    \label{eq:mu-binary}  \\
  \mu_3 &= \tfrac{1}{2} \qquad\text{(ternary)},
    \label{eq:mu-ternary}  \\
  \sigma_3 &= \tfrac{3}{2}.
    \label{eq:sigma-ternary}
\end{align}
\end{corollary}

\begin{proof}
$\mu_2 = 2 - 2/2 = 1$; $\mu_3 = 2 - 3/2 = 1/2$;
$\sigma_3 = 3/2$.  Each follows from
\cref{def:spectral-params} at the stated arity.
Verified by \lean{mu_binary}, \lean{mu_ternary},
\lean{sigma_ternary} in Lean.
\end{proof}


%%% --- Norm obstruction ---

\begin{lemma}[Diagonal contraction bound]\label{lem:diagonal-contraction}
If $0<t<1$, then $t^2 < t$. In particular, $t\le t^2$ is impossible.
\end{lemma}

\begin{proof}
$t - t^2 = t(1-t)$.  Since $0<t<1$, both $t>0$ and $1-t>0$, so
$t-t^2>0$, i.e., $t^2 < t$.
\end{proof}

\begin{theorem}[No-Diagonal Theorem]\label{thm:no-diagonal}
Let $\Omega$ be a tripartite object
(\cref{def:object-norm}) with
$\norm{\Omega}<1$. Then there is no evaluation
$\mathrm{eval}\colon\Omega\otimes\Omega\to\Omega$ with
$\mathrm{eval}\circ\Delta=\mathrm{id}_\Omega$ compatible with the norm
structure.
\end{theorem}

\begin{proof}
Suppose $\Delta\colon\Omega\to\Omega\otimes\Omega$ and
$\mathrm{eval}\colon\Omega\otimes\Omega\to\Omega$ exist with
$\mathrm{eval}\circ\Delta=\mathrm{id}_\Omega$, both compatible with the
norm structure:
%
\begin{center}
\begin{tikzcd}[column sep=4em]
  \Omega
    \arrow[r, "\Delta"]
    \arrow[rr, bend left=30, "\mathrm{id}_\Omega"]
  & \Omega \otimes \Omega
    \arrow[r, "\mathrm{eval}"]
  & \Omega
\end{tikzcd}
\end{center}
%
\noindent
By multiplicativity~(N2),
$\norm{\Omega\otimes\Omega}=\norm{\Omega}^2$.  For $\mathrm{eval}$ to
be norm-compatible, $\norm{\mathrm{eval}(X)}\le\norm{X}$
(submultiplicativity).  Tracking the norm through the diagram:
\[
  \norm{\Omega}
  \;\xrightarrow{\;\Delta\;}
  \norm{\Omega\otimes\Omega} = \norm{\Omega}^2
  \;\xrightarrow{\;\mathrm{eval}\;}
  \norm{\Omega},
\]
hence
$\norm{\Omega}=\norm{\mathrm{eval}\circ\Delta(\Omega)}\le
\norm{\Omega\otimes\Omega}=\norm{\Omega}^2$,
i.e., $\norm{\Omega}\le\norm{\Omega}^2$.  But $0<\norm{\Omega}<1$
contradicts \cref{lem:diagonal-contraction}.
\end{proof}

\begin{proposition}[Equivalence of obstructions]\label{prop:equiv-obstructions}
For $0 < \norm{\Omega} < 1$, the following are equivalent:
\\
(a) No exponential: $\Omega^\Omega$ is not representable.
\\
(b) No predominance: no point-surjection $\Omega\to\Omega^\Omega$.
\\
(c) No retraction: $\mathrm{eval}\circ\Delta\ne\mathrm{id}_\Omega$.
\\
All reduce to $\norm{\Omega}\le\norm{\Omega}^2$, which fails when $\norm{\Omega}<1$.
\end{proposition}

\begin{proof}
(a)$\Rightarrow$(b): If $\Omega^\Omega$ is not representable, no morphism
$\Omega\to\Omega^\Omega$ can exist.
(b)$\Rightarrow$(c): A retraction $\mathrm{eval}\circ\Delta=\mathrm{id}$
would compose with currying to give a predominance morphism.
(c)$\Rightarrow$(a): If $\Omega^\Omega$ were representable with evaluation
$\mathrm{eval}$, the diagonal $\Delta$ would give a retraction.  All three
conditions require $\norm{\Omega}\le\norm{\Omega}^2$, which fails by
\cref{lem:diagonal-contraction}.
\end{proof}

\begin{remark}[Exponential objects and cartesian closure]%
\label{rmk:exponential-objects}
The argument in \cref{prop:equiv-obstructions} is
\emph{conditional}: it does not assume that the ambient category is
cartesian closed, nor does it construct the exponential object
$\Omega^\Omega$.  Rather, it proves that $\Omega^\Omega$ \emph{cannot}
be representable when $\norm{\Omega}<1$.  This is precisely Lawvere's
pattern~\cite{lawvere69}: non-representability of the exponential is the
\textbf{conclusion}, not a hypothesis.  In categorical terms, the
subcategory of objects with $\norm{\cdot}<1$ fails to be cartesian
closed---and this failure is what blocks the diagonal arguments underlying
Russell, Cantor, and G\"odel.  No verification of cartesian closure is
required; the norm obstruction suffices.
\end{remark}

\begin{corollary}[Paradox avoidance]\label{cor:paradox-avoidance}
Tripartite objects with $\norm{\Omega}<1$ avoid Russell, Cantor,
G\"odel, and halting-problem paradoxes.
\end{corollary}

\begin{proof}
By Lawvere~\cite{lawvere69} and Yanofsky~\cite{yanofsky03}, each paradox
requires a diagonal morphism $\Delta$ with evaluation section.
\cref{prop:equiv-obstructions}(c) blocks all such sections when
$\norm{\Omega}<1$.
\end{proof}

%%% --- Arity Transition ---

\begin{proposition}[$\Set$: binary level]\label{prop:set-level}
$\mu=1$, diagonal permitted, Lawvere applies.
\end{proposition}

\begin{proof}
In $\Set$ with $\norm{X}=|X|$ and $\otimes=\times$, the effective arity is
$n=2$: $\mu_2=1$.  For $|X|\ge 2$: $|X|^{|X|}>|X|$ since
$|X|=|X|^1<|X|^{|X|}$.  The diagonal $\Delta(x)=(x,x)$ is permitted, and
$\norm{\Omega}=1$ does not trigger the contraction bound
($1\le 1^2 = 1$ holds).
\end{proof}

\begin{proposition}[$\Hilb$: the bridge]\label{prop:hilb-level}
Both algebraic norm ($\norm{I}=1$) and element norm
($\norm{\Omega}=\tfrac{1}{2}$, from~\cref{eq:norm-Omega-half}) coexist.
The diagonal is quadratic, not linear.
\end{proposition}

\begin{proof}
In $\Hilb$ with $\norm{H}=\dim(H)$ and tensor product, the unit has
$\norm{I}=1$ (algebraic).  The PCF element norm gives
$\norm{\Omega}=1/2$ (metric).  The diagonal map
$\Delta(v)=v\otimes v$ satisfies
$\Delta(\lambda v) = \lambda^2(v\otimes v)\ne\lambda\Delta(v)$ when
$\lambda\ne 0,1$: it is quadratic, not linear.
\end{proof}

\begin{proposition}[$\Met$: ternary level]\label{prop:met-level}
$\mu=1/2<1$
(\cref{eq:mu-ternary}), Lawvere is blocked.
\end{proposition}

\begin{proof}
In $\Met$ with $\norm{X}=\mathrm{diam}(X)$ and product metric, the
canonical norms $\norm{X_1}=1/\sqrt{3}$, $\norm{X_2}=1$,
$\norm{X_3}=\sqrt{3}/2$ give
$\norm{\Omega}=1/2<1$.  \cref{lem:diagonal-contraction}:
$1/2\le(1/2)^2=1/4$ fails.  Lawvere is blocked.
\end{proof}

\begin{theorem}[Arity Transition]\label{thm:arity-transition}
$\Set$, $\Hilb$, $\Met$ are levels of a single binary-to-ternary
transition connected by the forgetful functor.
\end{theorem}

\begin{proof}
The forgetful functor chain $U\colon C^*\text{--}\mathbf{Alg} \to \Hilb \to \Met \to \Set$ implements the transitions between the binary and ternary regimes:
%
\begin{center}
\begin{tikzcd}[column sep=3.5em, row sep=1.8em]
  \underset{\scriptscriptstyle \text{Binary arity}}{C^*\text{--}\mathbf{Alg}}
    \arrow[r, "U_1"]
  & \mathbf{Hilb}
    \arrow[r, "U_2"]
  & \underset{\scriptscriptstyle \text{Ternary arity}}{\mathbf{Met}}
    \arrow[r, "U_3"]
  & \mathbf{Set}
\end{tikzcd}
\end{center}
%
\noindent
The contraction factor is $\mu_2/\mu_3 = 1/(1/2) = 2 = |\ker G|$. The geometric capacity to support self-reference changes at each level:
\begin{itemize}
  \item \textbf{At $\Set$ ($n=2$):} $\mu=1$; arity-independent discrete structure. No metric restriction blocks the diagonal morphism (\cref{prop:set-level}).
  \item \textbf{At $\Hilb$:} Transition level. Norm coexistence of the algebraic unit ($\norm{I}_{\Hilb}=1$) and the metric PCF object ($\norm{\Omega}_{\text{\textsc{pcf}}} = 1/2$; \cref{eq:norm-Omega-half}). The diagonal is quadratic (\cref{prop:hilb-level}).
  \item \textbf{At $\Met$ ($n=3$):} $\mu=1/2 < 1$; arity-unique obstruction. Lawvere's diagonal is blocked by the contractive bound (\cref{prop:met-level,thm:no-diagonal}).
\end{itemize}
\end{proof}

%%% --- Spectral uniqueness ---

\begin{proposition}[Spectral uniqueness]\label{prop:spectral-uniqueness}
The system
\begin{equation}\label{eq:spectral-system}
  \sigma + \mu = 2, \qquad \sigma\,\mu = \tfrac{3}{4}
\end{equation}
has the unique solution $\sigma = 3/2$, $\mu = 1/2$.
\end{proposition}

\begin{proof}
Substituting $\mu=2-\sigma$ into $\sigma\mu=3/4$:
$\sigma(2-\sigma)=3/4$, i.e., $\sigma^2-2\sigma+3/4=0$.
The quadratic formula gives $\sigma=(2\pm\sqrt{4-3})/2=(2\pm 1)/2$.
Thus $\sigma\in\{3/2,\,1/2\}$.  Only $\sigma=3/2$ gives
$\mu=1/2<1$ (obstruction effective); $\sigma=1/2$ gives $\mu=3/2>1$,
violating the obstruction condition.
\end{proof}

%%% --- Two-level norms ---

\begin{definition}[Two-level norms]\label{def:two-level-norms}
The forgetful functor $U\colon C^*\text{--}\mathbf{Alg}\to\Hilb$ carries two coexisting norm levels:
\begin{enumerate}[label=\alph*), leftmargin=2em]
  \item \textbf{Algebra unit norm:} $\norm{\cdot}_{\text{\textsc{alg}}} = 1$ ($\Set$ level, $\mu_2 = 1$).
  \item \textbf{Element PCF norm:} $\norm{\Omega}_{\text{\textsc{pcf}}} = 1/2$ ($\Met$ level, $\mu_3 = 1/2$).
\end{enumerate}
The strict inequality $\norm{\Omega}_{\text{\textsc{pcf}}} < \norm{\cdot}_{\text{\textsc{alg}}}$ is the contraction that blocks the diagonal; the contraction factor is $1/2 = |\ker G|^{-1}$.
\end{definition}

%%% --- Galois functor ---

\begin{definition}[Galois functor]\label{def:galois-functor}
The \emph{Galois functor} is the surjective group homomorphism
\begin{equation}\label{eq:galois-functor}
  G\colon\Ztwenty\longrightarrow\mathbb{Z}_2\times\mathbb{Z}_2,
  \qquad
  G(p\bmod 20) =
  \bigl(\chi_4(p),\;\chi_5(p)\bigr),
\end{equation}
where $\chi_4 = \bigl(\frac{\cdot}{4}\bigr)$ and
$\chi_5 = \bigl(\frac{\cdot}{5}\bigr)$ are the Legendre symbols.
The four fibers partition $\Ztwenty$:
\begin{center}
\begin{tabular}{ccl}
  $G^{-1}(0,0)=\{1,9\}$, &\quad
  $G^{-1}(1,0)=\{11,19\}$, &\quad
  $G^{-1}(0,1)=\{13,17\}$, \\[3pt]
  \multicolumn{3}{c}{$G^{-1}(1,1)=\{3,7\}$.}
\end{tabular}
\end{center}
Verified by \lean{G_fibers_partition} (\texttt{decide}) in Lean.
\end{definition}

\begin{remark}[Categorical interpretation of $G$]%
\label{rmk:galois-categorical}
A group homomorphism $G\colon H\to K$ \emph{is} a functor between the
corresponding one-object categories: the single object~$\bullet$
maps to~$\bullet$, and each group element (viewed as a morphism
$\bullet\to\bullet$) maps to its image under~$G$.  Preservation of
composition is exactly the homomorphism property $G(ab)=G(a)G(b)$,
and preservation of the identity is $G(1)=1$.  For the Galois functor:
%
\begin{center}
\begin{tikzcd}[column sep=6em]
  \underset{\scriptstyle 8\text{ morphisms}}{\bullet}
    \arrow[r, "G"]
  & \underset{\scriptstyle 4\text{ morphisms}}{\bullet}
\end{tikzcd}
\end{center}
%
\noindent
The source category $\mathbf{B}\Ztwenty$ has $8$~morphisms (the
elements of~$\Ztwenty$); the target $\mathbf{B}(\mathbb{Z}_2\times
\mathbb{Z}_2)$ has $4$.  Surjectivity of~$G$ means the functor is
\emph{essentially surjective on morphisms}.  The kernel $\ker G =
\{1,9\}$ with $|\ker G|=2$ gives the contraction factor
$\mu_2/\mu_3=2$ (\cref{prop:contraction-factor}).
In the Lean formalization, $G$ is verified as a group homomorphism via
\texttt{decide}; the functorial interpretation follows by the standard
equivalence between groups and one-object categories.
\end{remark}

\begin{theorem}[Zeta factorization]\label{thm:zeta-factorization}
The Euler product of $\zeta(s)$ maps into the categorical framework
through four real-character $L$-functions, via
\cref{thm:zeta-forced} and $G$.
\end{theorem}

\begin{proof}
By \cref{thm:zeta-forced}, every prime $p>5$ maps to $\Ztwenty$
under reduction modulo~$20$.  The Galois functor
$G\colon\Ztwenty\to\mathbb{Z}_2\times\mathbb{Z}_2$
(\cref{thm:golden-prime-classes}) assigns to each factor
$(1-p^{-s})^{-1}$ of the Euler product a character value in
$\mathbb{Z}_2\times\mathbb{Z}_2$.  Since the group has exponent~$2$
(\cref{prop:exponent-two}), all four characters are real and
self-dual, yielding four $L$-functions $L(s,\tilde{\chi})$ with
$\tilde{\chi} = \chi \circ G$.
\end{proof}



% ------------------------------------------------------------------
\subsection{Representability Tower}\label{ssec:cat-limit}
% ------------------------------------------------------------------

The tower $\norm{\Omega}_k = (1/2)^{2^k}$ demonstrates that the No-Diagonal
obstruction is not an isolated phenomenon but persists \emph{at every level}
of the representability tower, growing exponentially stronger.  The tower
satisfies eight properties, preserved at every level~$k$, that collectively
constitute the tower convergence theorem.%
\footnote{The tower limit $\norm{\Omega}_k\to 0$ is a \emph{topological} 
limit of a real sequence (verified by \lean{Omega_norm_tendsto_zero} 
via \texttt{Filter.Tendsto}), not a categorical limit in the sense of 
a universal cone over a diagram. We use the term ``representability 
tower'' to avoid conflation with the categorical notion.}

\begin{definition}[Iterated norm tower]\label{def:Omega-norm-tower}
\begin{equation}\label{eq:Omega-norm}
  \norm{\Omega}_k = \left(\tfrac{1}{2}\right)^{2^k},
  \qquad k = 0,1,2,\ldots
\end{equation}
\end{definition}

\begin{lemma}[Tower properties]\label{lem:tower-properties}
The tower norm $\norm{\Omega}_k = (1/2)^{2^k}$ consists of eight properties (\textsc{l}1)--(\textsc{l}8) that assemble into the convergence structure. The first four define the sequence behavior for all $k \ge 0$:
\begin{enumerate}[label=(\textsc{l}\arabic*), leftmargin=3em]
  \item $0 < \norm{\Omega}_k < 1$ (Convergence bounds).
  \item $\norm{\Omega}_k < 1 \implies \Delta_k$ is blocked (Lawvere level check).
  \item $\norm{\Omega}_k \to 0$ as $k \to \infty$ (Topological limit).
  \item $\norm{\Omega}_{k+1} < \norm{\Omega}_k$ (Strictly decreasing sequence).
\end{enumerate}
Properties (\textsc{l}5)--(\textsc{l}8) integrate the structural and spectral witnesses: (\textsc{l}5) $\norm{\Omega}_0 = 1/2$; (\textsc{l}6) spectral uniqueness; (\textsc{l}7) additive destruction; and (\textsc{l}8) prime-forced ternary fragmentation.
\end{lemma}

\begin{proof}
(i) $\norm{\Omega}_k = (1/2)^{2^k} > 0$ since it is a positive power of
a positive number.
(ii) $(1/2)^{2^k} < 1$ since $1/2 < 1$ and $2^k \ge 1$.
(iii) $\norm{\Omega}_{k+1} = (1/2)^{2^{k+1}} = ((1/2)^{2^k})^2 =
\norm{\Omega}_k^2 < \norm{\Omega}_k$ since $0 < \norm{\Omega}_k < 1$.
(iv) $(1/2)^{2^k} \to 0$ as $k\to\infty$ since $2^k\to\infty$ and
$0 < 1/2 < 1$.
Verified by \lean{Omega_norm_pos}, \lean{Omega_norm_lt_one},
\lean{Omega_norm_decreasing}, \lean{Omega_norm_tendsto_zero} in Lean.
\end{proof}


\begin{corollary}[No-Diagonal at every level]\label{cor:no-diagonal-tower}
For all $k$, no diagonal morphism exists at level~$k$ (by
\cref{thm:no-diagonal} applied to $\norm{\Omega}_k < 1$).
No retraction is possible at any level.
\end{corollary}

\begin{proof}
By \cref{lem:tower-properties}, $0 < \norm{\Omega}_k < 1$ for
all~$k$.  By \cref{thm:no-diagonal}, $\norm{\Omega}_k < 1$
implies no evaluation-compatible diagonal exists at level~$k$.
\end{proof}


\begin{proposition}[Tower base]\label{prop:tower-base}
\begin{equation}\label{eq:tower-base}
  \norm{\Omega}_0 = \mu_3 = \tfrac{1}{2}
  \qquad\text{(\cref{eq:mu-ternary})}.
\end{equation}
\end{proposition}

\begin{proof}
$\norm{\Omega}_0 = (1/2)^{2^0} = (1/2)^1 = 1/2 = \mu_3$
(\cref{eq:mu-ternary}).
Verified by \lean{tower_base_is_mu3} in Lean.
\end{proof}


\begin{proposition}[Contraction factor]\label{prop:contraction-factor}
\begin{equation}\label{eq:contraction-factor}
  \frac{\mu_2}{\mu_3} = \frac{1}{1/2} = 2,
  \qquad\text{and}\quad
  \mu_3 < \mu_2 \quad\text{(strict contraction)}.
\end{equation}
\end{proposition}

\begin{proof}
$\mu_2/\mu_3 = 1/(1/2) = 2$.  $\mu_3 = 1/2 < 1 = \mu_2$.
\end{proof}


\begin{theorem}[Tower convergence theorem]\label{thm:cat-limit}
The tower $\{\norm{\Omega}_k\}$ converges to zero, with $\norm{\Omega}_0 =
\mu_3$. Components~(L1)--(L8) assemble into the representability tower
structure. The triple equality holds:
\begin{equation}\label{eq:triple-equality}
  \mu_3 = \norm{\Omega}_0 = \tfrac{1}{2}.
\end{equation}
\end{theorem}

\begin{proof}
Properties (L1)--(L3) and (L5) are \cref{lem:tower-properties}
and \cref{prop:tower-base}.  (L4) is the strict decrease from
\cref{lem:tower-properties}.  (L6) is
\cref{prop:spectral-uniqueness}.  (L7) is
\cref{prop:add-destroyed}.  (L8) is
\cref{thm:zeta-forced}.  The triple equality
$\mu_3 = \norm{\Omega}_0 = 1/2$ follows from (L5) and
\cref{eq:mu-ternary}.
Verified by \lean{categorical_limit_theorem} in Lean.
\end{proof}



% ------------------------------------------------------------------
\subsection{Omega Spectrum}\label{ssec:omega-spectrum}
% ------------------------------------------------------------------

The eigenvalues of $\OmH$ were defined in~\cref{eq:Omega-hat-eigenvalues}.
We now extract their spectral properties and reformulate the Riemann
Hypothesis in the language of the $\OmH$ operator.  The key observation is
that the three eigenvalues $\tfrac{1}{2},\,-\tfrac{1}{4}\pm\tfrac{\sqrt{3}}{4}i$
form an equilateral triangle inscribed in $|z|=1/2$, and only one of them
has positive real part---forcing any zero whose real part is an eigenvalue to
lie on the critical line.

\begin{proposition}[Properties of $\omega$]\label{prop:omega-properties}
The primitive cube root $\omega = e^{2\pi i/3}$
(\cref{eq:omega-root}) satisfies:
\begin{equation}\label{eq:omega-re-im}
  \mathrm{Re}(\omega) = -\tfrac{1}{2}, \qquad
  \mathrm{Im}(\omega) = \tfrac{\sqrt{3}}{2}.
\end{equation}
\end{proposition}

\begin{proof}
$\omega = e^{2\pi i/3} = \cos(2\pi/3) + i\sin(2\pi/3)$.
$\cos(2\pi/3) = -1/2$ and $\sin(2\pi/3) = \sqrt{3}/2$.
Verified by \lean{w_properties} in Lean.
\end{proof}


\begin{proposition}[Unique positive real part]\label{prop:Omega-hat-re}
\begin{equation}\label{eq:Omega-hat-0-re}
  \mathrm{Re}\bigl(\OmH(0)\bigr) = \tfrac{1}{2},
\end{equation}
and this is the \emph{unique} eigenvalue with positive real part:
\begin{equation}\label{eq:Omega-hat-unique-positive-re}
  \mathrm{Re}\bigl(\OmH(k)\bigr) > 0 \implies
  \mathrm{Re}\bigl(\OmH(k)\bigr) = \tfrac{1}{2}.
\end{equation}
\end{proposition}

\begin{proof}
$\widehat{\Omega}(0) = \tfrac{1}{2}\omega^0 = 1/2$, so
$\mathrm{Re}(\widehat{\Omega}(0)) = 1/2 > 0$.
For $k=1$: $\mathrm{Re}(\widehat{\Omega}(1)) = \mathrm{Re}(\tfrac{1}{2}\omega) =
\tfrac{1}{2}\cdot(-\tfrac{1}{2}) = -1/4 < 0$.
For $k=2$: $\mathrm{Re}(\widehat{\Omega}(2)) = -1/4 < 0$ by conjugation.
Hence $k=0$ is the unique index with positive real part, and
$\mathrm{Re}(\widehat{\Omega}(0)) = 1/2$.
Verified by \lean{Omega_hat_unique_positive_re} in Lean.
\end{proof}


\begin{proposition}[Spectral embedding]\label{prop:spectral-embedding}
Let $\rho$ be a non-trivial zero of~$\zeta(s)$ with
$0<\mathrm{Re}(\rho)<1$.  Within the PCF categorical framework:
\begin{equation}\label{eq:spectral-embedding}
  \mathrm{Re}(\rho) \;\in\;
  \mathrm{Re}\bigl(\mathrm{spec}(\OmH)\bigr)
  \;=\; \bigl\{\tfrac{1}{2},\; -\tfrac{1}{4}\bigr\}.
\end{equation}
\end{proposition}

\begin{proof}
The chain proceeds in four steps.

\emph{Step~1: Euler product enters~$\Ztwenty$.}\label{step:1}
The Euler product $\zeta(s)=\prod_p(1-p^{-s})^{-1}$ is a product over
primes.  Every prime $p>5$ satisfies $p\bmod 20\in\Ztwenty$
(\cref{thm:zeta-forced}).  Therefore each factor of the Euler
product is indexed by an element of~$\Ztwenty$.

\emph{Step~2: $\Ztwenty$ forces the ternary fragment.}\label{step:2}
The group $\Ztwenty$ is closed under multiplication but not under
addition (\cref{prop:add-destroyed}).  This purely
multiplicative arithmetic selects the ternary fragment with
$\mu_3=1/2$ (\cref{prop:fragment-mu}), uniquely determined
by the quadratic system $\sigma+\mu=2$, $\sigma\mu=3/4$
(\cref{prop:spectral-uniqueness}).

\emph{Step~3: The ternary fragment has operator~$\OmH$.}\label{step:3}
The ternary fragment is generated by the tripartite object
$\Omega = P\cdot C\cdot F$ with $\norm{\Omega}=1/2$
(\cref{eq:norm-Omega-half}), whose spectral decomposition is
$\OmH = \tfrac{1}{2}\,\mathrm{diag}(1,\omega,\omega^2)$
(\cref{def:Omega-hat}).  The No-Diagonal
\cref{thm:no-diagonal} blocks self-reference at
$\norm{\Omega}<1$; the Arity Transition
(\cref{thm:arity-transition}) identifies this contraction as
the forgetful functor $U\colon C^*$-$\mathbf{Alg}\to\Hilb\to\Met$
carrying the norm from~$1$ to~$1/2$.

\emph{Step~4: Zeros inherit the spectral constraint.}\label{step:4}
The multiplicative structure of~$\Ztwenty$ groups Euler factors into
exactly four classes via the Galois functor (\cref{def:galois-functor}),
producing four real Dirichlet characters (all of which have
exponent~$2$, \cref{prop:exponent-two}).
The tripartite $S_3$
symmetry of the generating object $\Omega = P\cdot C\cdot F$ constrains
any norm-compatible operator to have spectrum
$\mathrm{spec}(\OmH) = \tfrac{1}{2}\{1,\omega,\omega^2\}$, whose real
parts are $\{1/2, -1/4\}$ (\cref{prop:Omega-hat-re}).
Since each $L(s,\tilde\chi_j)$ factorizes through primes in~$\Ztwenty$
(\cref{step:1}) and the fragment's spectral parameter is $\mu_3 = 1/2$ (\cref{step:2}),
a zero~$\rho$ in the critical strip satisfies
$\mathrm{Re}(\rho) \in \mathrm{Re}(\mathrm{spec}(\OmH))$.
\emph{Note:} This step combines the Lean-verified arithmetic
(\cref{step:1,step:2,step:3}) with the classical spectral mapping
theorem
(Dunford--Schwartz~\cite{dunford58}) and the Euler product identity.
The Lean formalization in \texttt{PCF\_Emulation.lean} makes this
dependence explicit via two named axioms (\texttt{bridge\_axiom},
\texttt{hecke\_companion}) whose classical justifications are
Gelfand~(1941)~\cite{gelfand41}, Dunford--Schwartz~\cite{dunford58},
and Hecke~(1920)~\cite{hecke20}.
The concrete structure on which these classical tools act---the
arithmetic--spectral span with common source~$G$---is constructed
in~\cite{gonzalez26}.
\end{proof}


\begin{theorem}[Riemann Hypothesis --- spectral form]\label{thm:rh-spectral}
Let $\rho$ be a non-trivial zero of~$\zeta(s)$ with
$0<\mathrm{Re}(\rho)<1$.  Within the PCF categorical framework,
$\mathrm{Re}(\rho) = 1/2$.
\end{theorem}

\begin{proof}[First proof (spectral)]
By \cref{prop:spectral-embedding},
$\mathrm{Re}(\rho)\in\{1/2,\,-1/4\}$.  Since
$\mathrm{Re}(\rho)>0$, the value $-1/4$ is excluded.  Hence
$\mathrm{Re}(\rho)=1/2$.
\end{proof}


\begin{theorem}[Operator $\omega$ master theorem]\label{thm:omega-master}
The operator $\OmH$ encodes five properties simultaneously:
\begin{enumerate}[label=(\arabic*), leftmargin=2.5em]
  \item $\OmH(k) = \tfrac{1}{2}\omega^k$ (eigenvalue formula, \cref{eq:Omega-hat-eigenvalues}).
  \item $|\OmH(k)| = 1/2$ for all $k$ (constant modulus, \cref{eq:modulus-Omega}).
  \item $\mathrm{Re}(\OmH(0)) = 1/2$ (critical line, \cref{eq:Omega-hat-0-re}).
  \item Unique positive-Re eigenvalue (\cref{eq:Omega-hat-unique-positive-re}).
  \item $S_3$ symmetry: the three eigenvalues are related by $120^\circ$ rotations.
\end{enumerate}
\end{theorem}

\begin{proof}
Properties (1)--(4) follow from the definitions
(\cref{eq:Omega-hat-eigenvalues,eq:modulus-Omega,eq:Omega-hat-0-re,eq:Omega-hat-unique-positive-re}).
Property~(5): the three eigenvalues $\tfrac{1}{2}\omega^k$ for $k=0,1,2$
are related by multiplication by $\omega = e^{2\pi i/3}$, i.e.,
$120^\circ$ rotations.
\end{proof}



% ------------------------------------------------------------------
